% Données : sources, collecte, gel, référentiel, qualité, scellé (parties 1 et 2 ; DATA_FREEZE.md).

\chapter{Données}
\label{chap:donnees}

\section{Sources}

\begin{center}
\begin{tabularx}{\textwidth}{lXl}
\toprule
Source & Contenu utilisé & Statut \\
\midrule
API-FOOTBALL & matchs, compositions, statistiques d'équipe et de joueurs, profils, entraîneurs,
               classements, depuis 2015-16 & abonnement terminé le 22/10/2026 \\
football-data & résultats depuis 2000, tirs, cotes plus ou moins 2,5 buts ; résultats et prochains
                matchs en direct & gratuit, deux mises à jour par semaine \\
\bottomrule
\end{tabularx}
\end{center}

Understat (xG) est écarté : son \code{robots.txt} interdit la collecte automatique (ADR-0023). Les xG
d'API-FOOTBALL ne couvrant que les saisons récentes, le projet construit un \code{xg\_proxy} à partir
des tirs cadrés (ADR-0032, ADR-0035).

\section{Collecte et gel}

\paragraph{Brut d'abord.} Chaque réponse de l'API est écrite telle quelle dans un fichier
\code{.json.gz}, avec une ligne dans un journal de requêtes (ADR-0003) ; les CSV de football-data sont
gardés octet pour octet (ADR-0024). Rien n'est jamais modifié dans ce dossier : tout le reste se
reconstruit à partir de lui.

\paragraph{Collecteur v2.} Le collecteur demande les matchs par lots de 20 (un seul appel rend
compositions, événements et statistiques), lit le champ d'erreurs de l'API, reprend après une
interruption et prend un verrou avant toute écriture (ADR-0004, ADR-0018). Le périmètre est collecté
par paliers : top 5 et deuxièmes divisions de 2015-16 à 2026-27, coupes d'Europe et coupes nationales,
autres championnats, classements et profils de joueurs (ADR-0002).

\paragraph{Gel du 19 octobre 2026.} Avant la fin de l'abonnement, le brut est figé (tag
\code{data-freeze-2026-10}) : \textbf{17\,336 fichiers}, 306~Mo compressés, \textbf{102\,554 matchs
terminés}, tous avec leur détail. Le brut est sauvegardé sur un disque externe, puis restauré et
recompté pour vérifier la sauvegarde (ADR-0006). Un journal \og T-60 \fg{} a aussi enregistré, du 9 au
18 octobre, les compositions publiées une heure avant le coup d'envoi, pour la version
intermédiaire (ADR-0017).

\section{Référentiel}

La commande \code{load} reconstruit entièrement le schéma \code{staging} de Postgres à partir du brut :
compétitions, équipes, joueurs, matchs, compositions, tirs et cotes (ADR-0008). Les identifiants sont
ceux d'API-FOOTBALL, jamais les noms : deux joueurs homonymes restent distincts. Les corrections
(équipes de football-data à rapprocher, collisions d'identifiants) sont des fichiers YAML relus par
\code{load}, jamais des modifications à la main dans la base. Sur le brut définitif, \code{load}
produit 167\,980 matchs en un peu plus de 8 minutes. Deux chargements successifs donnent 15 empreintes
de tables sur 15 identiques : la reconstruction est reproductible.

\section{Qualité}

Des contrôles chiffrés accompagnent chaque chargement (\code{check-referentiel}) :

\begin{itemize}
  \item appariement des matchs de football-data et d'API-FOOTBALL : 99,98~\% (seuil : 99,5~\%) ;
  \item identité des joueurs : collisions d'identifiants exclues automatiquement et comptées
        (ADR-0020) ; 0,14~\% des matchs du top 5 touchés par une exclusion (seuil : 1~\%) ;
  \item ruptures de série signalées plutôt que corrigées : les tirs de football-data de la Serie A
        de 2018-19 à 2020-21 ne se comparent pas aux autres saisons (ADR-0029), d'où une source de tirs
        choisie saison par saison (ADR-0035).
\end{itemize}

\section{Le scellé}

Tous les matchs joués à partir du 1\textsuperscript{er} juillet 2025 sont \emph{sous scellés}
(ADR-0012) : aucune de leurs valeurs (score, tirs, cotes) n'a été lue pendant le développement. Ce
n'est pas qu'une consigne : toute lecture des matchs passe par une seule fonction,
\code{features/sources.py}, qui filtre la date en SQL, et des tests d'architecture vérifient
qu'aucun autre module n'interroge la table des matchs (ADR-0028). Le scellé a servi à deux tests, un par
horizon, chacun une seule fois (chapitre~\ref{chap:resultats}). Depuis le premier, l'inférence seule peut
lire les matchs scellés, comme historique de variables (ADR-0045) ; depuis le second, elle peut aussi les
rejouer, avec leur score mais sans leurs cotes (ADR-0052). Le développement ne les lit toujours pas.

\section{Compositions et entraîneurs}

L'horizon H2 (chapitre~\ref{chap:methode}) s'appuie sur des données déjà sauvegardées au gel : pour chaque
match, la feuille de match (onze de départ, remplaçants, entraîneur, formation) et, pour chaque joueur
inscrit, ses minutes, sa note et son poste, avec le profil du joueur (date de naissance). Elles
viennent d'API-FOOTBALL depuis 2015-16 (entraîneurs depuis 2016-17) et sont lues par la même porte unique
que les matchs, qui filtre le scellé en SQL. Aucun \code{load} supplémentaire n'a été nécessaire. Deux
limites en découlent : les notes sont calculées après le match par le fournisseur, selon une formule
inconnue, et seuls les matchs \emph{antérieurs} servent ; et aucune composition n'arrive plus après la fin
de l'abonnement, d'où un horizon H2 de rejeu seulement.
